%% ============================================================================
%% Related Work — Checkpoint 1 (expanded).
%%
%% Uses ten author-verified references in references.bib:
%%   slater2016enhancing, wu2020effectiveness, tsivitanidou2021inquiry,
%%   georgiou2021led, makransky2019gender, wang2019haptic, lecuyer2000pseudo,
%%   ujitoko2021pseudo, kim2022multisensory, bowman2004
%%
%% To use: %% ============================================================================
%% Related Work — Checkpoint 1 (expanded).
%%
%% Uses ten author-verified references in references.bib:
%%   slater2016enhancing, wu2020effectiveness, tsivitanidou2021inquiry,
%%   georgiou2021led, makransky2019gender, wang2019haptic, lecuyer2000pseudo,
%%   ujitoko2021pseudo, kim2022multisensory, bowman2004
%%
%% To use: %% ============================================================================
%% Related Work — Checkpoint 1 (expanded).
%%
%% Uses ten author-verified references in references.bib:
%%   slater2016enhancing, wu2020effectiveness, tsivitanidou2021inquiry,
%%   georgiou2021led, makransky2019gender, wang2019haptic, lecuyer2000pseudo,
%%   ujitoko2021pseudo, kim2022multisensory, bowman2004
%%
%% To use: %% ============================================================================
%% Related Work — Checkpoint 1 (expanded).
%%
%% Uses ten author-verified references in references.bib:
%%   slater2016enhancing, wu2020effectiveness, tsivitanidou2021inquiry,
%%   georgiou2021led, makransky2019gender, wang2019haptic, lecuyer2000pseudo,
%%   ujitoko2021pseudo, kim2022multisensory, bowman2004
%%
%% To use: \input{related-work} in main.tex in place of the current
%% \section{Related Work}.
%% ============================================================================

\section{Related Work}
Our sandbox sits at the intersection of four lines of research: immersive
virtual reality (VR) for science learning, the difficulty of rendering weight and
force on commodity hardware, the pseudo-haptic techniques that stand in for that
missing sensation, and the interaction-design foundations that make direct
manipulation legible. We review each in turn and close by stating what
distinguishes our approach.

\subsection{Immersive VR for Science Learning}
The pedagogical case for immersive VR rests less on visual fidelity than on the
sense of being present. Slater and S{\'a}nchez-Vives frame this through the place
illusion (the feeling of being in the virtual space) and the plausibility
illusion (the feeling that events are really happening), which together lead
users to respond to virtual events as if they were real~\cite{slater2016enhancing}.
For a physics sandbox this matters directly. If a learner treats a virtual block
and lever as real objects, then manipulating them can carry the intuitive weight
that a 2D simulation cannot.

Empirical work in physics education adds a consistent caveat: immersion helps,
but only when the learning experience is designed for it. A recent meta-analysis
by Wu et al.\ across many head-mounted display studies finds that immersive VR
yields a positive but moderate overall effect on learning performance, and that
this effect is strongly moderated by factors such as subject domain, learner
control, and how the experience is designed~\cite{wu2020effectiveness}.
Tsivitanidou, Georgiou, and Ioannou report on an inquiry-based immersive VR
physics activity from two angles: the learning experience they observe conceptual
gains alongside an interaction with learners' attitudinal
profiles~\cite{tsivitanidou2021inquiry}, and the design of that same experience,
which they document as a worked case of learning-experience
design~\cite{georgiou2021led}. We read this pair as evidence that immersive VR
physics learning is feasible and can produce conceptual gains; it establishes
that a VR physics activity can work, though it does not itself isolate felt weight
or the effect of a haptic cue. One common way to scaffold such experiences is an
embedded pedagogical agent. Makransky et al.\ studied one such agent and found a
gender-matching effect: a virtual instructor aided learning most when its
apparent gender matched the learner, a specific interaction rather than a general
endorsement of pedagogical agents~\cite{makransky2019gender}. Taken together,
these studies establish that immersive physics learning is promising but
design-sensitive. The value comes from what the learner is asked to do and how
the task is scaffolded, not from the headset itself. We take a deliberately
different scaffolding stance. Rather than an in-scene agent, we scaffold through
scene and task design, using tightly scoped scenarios that isolate a single
variable so that the cause of each outcome is unambiguous.

\subsection{The Missing Channel: Rendering Weight and Force}
The central obstacle for a physics-of-materials experience is that the very
quantity the lessons depend on, namely weight and the force needed to move a
loaded mechanism, is the one commodity that VR cannot deliver. A handheld
controller weighs the same whatever it represents. Wang et al.\ survey the state
of haptic display for VR and make clear that genuine, general-purpose force
feedback remains far from consumer reach, constrained by actuation, workspace,
and grounding requirements that headset-and-controller setups do not
meet~\cite{wang2019haptic}. Their survey frames the design problem for us
precisely. If we want a denser object to feel heavier on a Quest 2, we cannot
wait for the hardware to render the force, and must instead substitute perception
for physics.

\subsection{Pseudo-Haptics: Substituting Perception for Physics}
Pseudo-haptics addresses this gap by exploiting the dominance of vision over
proprioception. Manipulating the visual consequences of a motion, such as the
control-to-display gain or a small lag between hand and object, induces a felt
sense of weight, stiffness, or resistance without any force actuator. The
foundational demonstration comes from L{\'e}cuyer et al., who showed that
altering the visual response of an isometric input device makes users perceive
force and stiffness that no actuator produces, establishing that a
control-to-display manipulation alone can stand in for force
feedback~\cite{lecuyer2000pseudo}. Ujitoko and Ban survey the design space and
applications of pseudo-haptics, cataloguing how such visual manipulations
reliably produce haptic sensations and how they have been
applied~\cite{ujitoko2021pseudo}. Their taxonomy is the methodological backbone
for how we map mass and mechanical load onto gain and visual lag.

Two results ground our specific cue, the felt \emph{weight} of a held object.
Dominjon et al.\ manipulated the control/display ratio of a manipulated object in
a virtual environment and showed that this ratio directly changes its perceived
mass, so that the same physical motion can be made to feel heavier or lighter by
scaling its visual displacement~\cite{dominjon2005influence}. Rietzler et al.\
carried the same idea to consumer VR: by introducing perceivable offsets between
the tracked hand and the rendered object, they elicited a reliable sensation of
weight on a commodity headset with no force
hardware~\cite{rietzler2018breaking}. Together these establish that a gain or
tracking-offset manipulation can render felt weight on exactly the class of
device we target, and they fix the design variable for us: the cue must track an
object's \emph{mass}, the quantity that determines weight, rather than its
density. These techniques are not limited to lifting a single object. The same
principle of modulating the visual outcome of an effort extends naturally to the
resistance of operating a machine, such as a lever that grows harder to move as
its load increases.

\subsection{Foundations of 3D Interaction}
Whatever the illusion, the learner must first be able to reach out, grab, and
place objects fluidly, or the physics never gets probed. We ground our
interaction design in Bowman et al.'s canonical treatment of 3D user interfaces,
which formalizes selection, manipulation, and travel as the core tasks of spatial
interaction and catalogues the techniques for each~\cite{bowman2004}. This
foundation informs the grab-and-place model at the heart of the sandbox and keeps
our contribution, the pseudo-haptic property mapping, layered on top of
well-understood manipulation techniques rather than reinventing them.

\subsection{What Is Different About Our Approach}
Two gaps emerge from this literature. First, immersive physics studies
consistently deliver visual and interactive experiences but rarely give learners
a genuine felt contrast in weight between objects. A student may see that two
blocks differ, but does not feel one resist more than the other. Second,
pseudo-haptic studies convincingly render the felt weight of an isolated
object~\cite{dominjon2005influence, rietzler2018breaking}, but seldom span both
material properties (density, buoyancy, friction) and simple machines (levers,
pulleys) within one coherent environment and interaction model. Our sandbox
targets this intersection. We use pseudo-haptic gain and visual-lag modulation as
a unifying mechanism that makes an unrenderable physical quantity perceptible
across a family of controlled, single-variable scenarios, and we extend the same
mechanism from holding a heavier object to feeling a loaded lever push back.
Crucially, the cue tracks \emph{mass} for held objects and \emph{load} for the
lever, the quantities that actually determine weight and effort, so that density
is taught where it belongs, through the volume difference between equal-mass
blocks and through float-or-sink behavior, rather than being conflated with
heaviness. Because our contribution is an interaction technique for making a
physical quantity felt, we evaluate it at two levels: a system level, checking
that the cues behave monotonically with the underlying physics and that each
scenario matches its closed-form prediction; and a human level, a within-subjects
study testing whether the weight cue changes how accurately learners predict
outcomes that depend on mass, density, and load.
 in main.tex in place of the current
%% \section{Related Work}.
%% ============================================================================

\section{Related Work}
Our sandbox sits at the intersection of four lines of research: immersive
virtual reality (VR) for science learning, the difficulty of rendering weight and
force on commodity hardware, the pseudo-haptic techniques that stand in for that
missing sensation, and the interaction-design foundations that make direct
manipulation legible. We review each in turn and close by stating what
distinguishes our approach.

\subsection{Immersive VR for Science Learning}
The pedagogical case for immersive VR rests less on visual fidelity than on the
sense of being present. Slater and S{\'a}nchez-Vives frame this through the place
illusion (the feeling of being in the virtual space) and the plausibility
illusion (the feeling that events are really happening), which together lead
users to respond to virtual events as if they were real~\cite{slater2016enhancing}.
For a physics sandbox this matters directly. If a learner treats a virtual block
and lever as real objects, then manipulating them can carry the intuitive weight
that a 2D simulation cannot.

Empirical work in physics education adds a consistent caveat: immersion helps,
but only when the learning experience is designed for it. A recent meta-analysis
by Wu et al.\ across many head-mounted display studies finds that immersive VR
yields a positive but moderate overall effect on learning performance, and that
this effect is strongly moderated by factors such as subject domain, learner
control, and how the experience is designed~\cite{wu2020effectiveness}.
Tsivitanidou, Georgiou, and Ioannou report on an inquiry-based immersive VR
physics activity from two angles: the learning experience they observe conceptual
gains alongside an interaction with learners' attitudinal
profiles~\cite{tsivitanidou2021inquiry}, and the design of that same experience,
which they document as a worked case of learning-experience
design~\cite{georgiou2021led}. We read this pair as evidence that immersive VR
physics learning is feasible and can produce conceptual gains; it establishes
that a VR physics activity can work, though it does not itself isolate felt weight
or the effect of a haptic cue. One common way to scaffold such experiences is an
embedded pedagogical agent. Makransky et al.\ studied one such agent and found a
gender-matching effect: a virtual instructor aided learning most when its
apparent gender matched the learner, a specific interaction rather than a general
endorsement of pedagogical agents~\cite{makransky2019gender}. Taken together,
these studies establish that immersive physics learning is promising but
design-sensitive. The value comes from what the learner is asked to do and how
the task is scaffolded, not from the headset itself. We take a deliberately
different scaffolding stance. Rather than an in-scene agent, we scaffold through
scene and task design, using tightly scoped scenarios that isolate a single
variable so that the cause of each outcome is unambiguous.

\subsection{The Missing Channel: Rendering Weight and Force}
The central obstacle for a physics-of-materials experience is that the very
quantity the lessons depend on, namely weight and the force needed to move a
loaded mechanism, is the one commodity that VR cannot deliver. A handheld
controller weighs the same whatever it represents. Wang et al.\ survey the state
of haptic display for VR and make clear that genuine, general-purpose force
feedback remains far from consumer reach, constrained by actuation, workspace,
and grounding requirements that headset-and-controller setups do not
meet~\cite{wang2019haptic}. Their survey frames the design problem for us
precisely. If we want a denser object to feel heavier on a Quest 2, we cannot
wait for the hardware to render the force, and must instead substitute perception
for physics.

\subsection{Pseudo-Haptics: Substituting Perception for Physics}
Pseudo-haptics addresses this gap by exploiting the dominance of vision over
proprioception. Manipulating the visual consequences of a motion, such as the
control-to-display gain or a small lag between hand and object, induces a felt
sense of weight, stiffness, or resistance without any force actuator. The
foundational demonstration comes from L{\'e}cuyer et al., who showed that
altering the visual response of an isometric input device makes users perceive
force and stiffness that no actuator produces, establishing that a
control-to-display manipulation alone can stand in for force
feedback~\cite{lecuyer2000pseudo}. Ujitoko and Ban survey the design space and
applications of pseudo-haptics, cataloguing how such visual manipulations
reliably produce haptic sensations and how they have been
applied~\cite{ujitoko2021pseudo}. Their taxonomy is the methodological backbone
for how we map mass and mechanical load onto gain and visual lag.

Two results ground our specific cue, the felt \emph{weight} of a held object.
Dominjon et al.\ manipulated the control/display ratio of a manipulated object in
a virtual environment and showed that this ratio directly changes its perceived
mass, so that the same physical motion can be made to feel heavier or lighter by
scaling its visual displacement~\cite{dominjon2005influence}. Rietzler et al.\
carried the same idea to consumer VR: by introducing perceivable offsets between
the tracked hand and the rendered object, they elicited a reliable sensation of
weight on a commodity headset with no force
hardware~\cite{rietzler2018breaking}. Together these establish that a gain or
tracking-offset manipulation can render felt weight on exactly the class of
device we target, and they fix the design variable for us: the cue must track an
object's \emph{mass}, the quantity that determines weight, rather than its
density. These techniques are not limited to lifting a single object. The same
principle of modulating the visual outcome of an effort extends naturally to the
resistance of operating a machine, such as a lever that grows harder to move as
its load increases.

\subsection{Foundations of 3D Interaction}
Whatever the illusion, the learner must first be able to reach out, grab, and
place objects fluidly, or the physics never gets probed. We ground our
interaction design in Bowman et al.'s canonical treatment of 3D user interfaces,
which formalizes selection, manipulation, and travel as the core tasks of spatial
interaction and catalogues the techniques for each~\cite{bowman2004}. This
foundation informs the grab-and-place model at the heart of the sandbox and keeps
our contribution, the pseudo-haptic property mapping, layered on top of
well-understood manipulation techniques rather than reinventing them.

\subsection{What Is Different About Our Approach}
Two gaps emerge from this literature. First, immersive physics studies
consistently deliver visual and interactive experiences but rarely give learners
a genuine felt contrast in weight between objects. A student may see that two
blocks differ, but does not feel one resist more than the other. Second,
pseudo-haptic studies convincingly render the felt weight of an isolated
object~\cite{dominjon2005influence, rietzler2018breaking}, but seldom span both
material properties (density, buoyancy, friction) and simple machines (levers,
pulleys) within one coherent environment and interaction model. Our sandbox
targets this intersection. We use pseudo-haptic gain and visual-lag modulation as
a unifying mechanism that makes an unrenderable physical quantity perceptible
across a family of controlled, single-variable scenarios, and we extend the same
mechanism from holding a heavier object to feeling a loaded lever push back.
Crucially, the cue tracks \emph{mass} for held objects and \emph{load} for the
lever, the quantities that actually determine weight and effort, so that density
is taught where it belongs, through the volume difference between equal-mass
blocks and through float-or-sink behavior, rather than being conflated with
heaviness. Because our contribution is an interaction technique for making a
physical quantity felt, we evaluate it at two levels: a system level, checking
that the cues behave monotonically with the underlying physics and that each
scenario matches its closed-form prediction; and a human level, a within-subjects
study testing whether the weight cue changes how accurately learners predict
outcomes that depend on mass, density, and load.
 in main.tex in place of the current
%% \section{Related Work}.
%% ============================================================================

\section{Related Work}
Our sandbox sits at the intersection of four lines of research: immersive
virtual reality (VR) for science learning, the difficulty of rendering weight and
force on commodity hardware, the pseudo-haptic techniques that stand in for that
missing sensation, and the interaction-design foundations that make direct
manipulation legible. We review each in turn and close by stating what
distinguishes our approach.

\subsection{Immersive VR for Science Learning}
The pedagogical case for immersive VR rests less on visual fidelity than on the
sense of being present. Slater and S{\'a}nchez-Vives frame this through the place
illusion (the feeling of being in the virtual space) and the plausibility
illusion (the feeling that events are really happening), which together lead
users to respond to virtual events as if they were real~\cite{slater2016enhancing}.
For a physics sandbox this matters directly. If a learner treats a virtual block
and lever as real objects, then manipulating them can carry the intuitive weight
that a 2D simulation cannot.

Empirical work in physics education adds a consistent caveat: immersion helps,
but only when the learning experience is designed for it. A recent meta-analysis
by Wu et al.\ across many head-mounted display studies finds that immersive VR
yields a positive but moderate overall effect on learning performance, and that
this effect is strongly moderated by factors such as subject domain, learner
control, and how the experience is designed~\cite{wu2020effectiveness}.
Tsivitanidou, Georgiou, and Ioannou report on an inquiry-based immersive VR
physics activity from two angles: the learning experience they observe conceptual
gains alongside an interaction with learners' attitudinal
profiles~\cite{tsivitanidou2021inquiry}, and the design of that same experience,
which they document as a worked case of learning-experience
design~\cite{georgiou2021led}. We read this pair as evidence that immersive VR
physics learning is feasible and can produce conceptual gains; it establishes
that a VR physics activity can work, though it does not itself isolate felt weight
or the effect of a haptic cue. One common way to scaffold such experiences is an
embedded pedagogical agent. Makransky et al.\ studied one such agent and found a
gender-matching effect: a virtual instructor aided learning most when its
apparent gender matched the learner, a specific interaction rather than a general
endorsement of pedagogical agents~\cite{makransky2019gender}. Taken together,
these studies establish that immersive physics learning is promising but
design-sensitive. The value comes from what the learner is asked to do and how
the task is scaffolded, not from the headset itself. We take a deliberately
different scaffolding stance. Rather than an in-scene agent, we scaffold through
scene and task design, using tightly scoped scenarios that isolate a single
variable so that the cause of each outcome is unambiguous.

\subsection{The Missing Channel: Rendering Weight and Force}
The central obstacle for a physics-of-materials experience is that the very
quantity the lessons depend on, namely weight and the force needed to move a
loaded mechanism, is the one commodity that VR cannot deliver. A handheld
controller weighs the same whatever it represents. Wang et al.\ survey the state
of haptic display for VR and make clear that genuine, general-purpose force
feedback remains far from consumer reach, constrained by actuation, workspace,
and grounding requirements that headset-and-controller setups do not
meet~\cite{wang2019haptic}. Their survey frames the design problem for us
precisely. If we want a denser object to feel heavier on a Quest 2, we cannot
wait for the hardware to render the force, and must instead substitute perception
for physics.

\subsection{Pseudo-Haptics: Substituting Perception for Physics}
Pseudo-haptics addresses this gap by exploiting the dominance of vision over
proprioception. Manipulating the visual consequences of a motion, such as the
control-to-display gain or a small lag between hand and object, induces a felt
sense of weight, stiffness, or resistance without any force actuator. The
foundational demonstration comes from L{\'e}cuyer et al., who showed that
altering the visual response of an isometric input device makes users perceive
force and stiffness that no actuator produces, establishing that a
control-to-display manipulation alone can stand in for force
feedback~\cite{lecuyer2000pseudo}. Ujitoko and Ban survey the design space and
applications of pseudo-haptics, cataloguing how such visual manipulations
reliably produce haptic sensations and how they have been
applied~\cite{ujitoko2021pseudo}. Their taxonomy is the methodological backbone
for how we map mass and mechanical load onto gain and visual lag.

Two results ground our specific cue, the felt \emph{weight} of a held object.
Dominjon et al.\ manipulated the control/display ratio of a manipulated object in
a virtual environment and showed that this ratio directly changes its perceived
mass, so that the same physical motion can be made to feel heavier or lighter by
scaling its visual displacement~\cite{dominjon2005influence}. Rietzler et al.\
carried the same idea to consumer VR: by introducing perceivable offsets between
the tracked hand and the rendered object, they elicited a reliable sensation of
weight on a commodity headset with no force
hardware~\cite{rietzler2018breaking}. Together these establish that a gain or
tracking-offset manipulation can render felt weight on exactly the class of
device we target, and they fix the design variable for us: the cue must track an
object's \emph{mass}, the quantity that determines weight, rather than its
density. These techniques are not limited to lifting a single object. The same
principle of modulating the visual outcome of an effort extends naturally to the
resistance of operating a machine, such as a lever that grows harder to move as
its load increases.

\subsection{Foundations of 3D Interaction}
Whatever the illusion, the learner must first be able to reach out, grab, and
place objects fluidly, or the physics never gets probed. We ground our
interaction design in Bowman et al.'s canonical treatment of 3D user interfaces,
which formalizes selection, manipulation, and travel as the core tasks of spatial
interaction and catalogues the techniques for each~\cite{bowman2004}. This
foundation informs the grab-and-place model at the heart of the sandbox and keeps
our contribution, the pseudo-haptic property mapping, layered on top of
well-understood manipulation techniques rather than reinventing them.

\subsection{What Is Different About Our Approach}
Two gaps emerge from this literature. First, immersive physics studies
consistently deliver visual and interactive experiences but rarely give learners
a genuine felt contrast in weight between objects. A student may see that two
blocks differ, but does not feel one resist more than the other. Second,
pseudo-haptic studies convincingly render the felt weight of an isolated
object~\cite{dominjon2005influence, rietzler2018breaking}, but seldom span both
material properties (density, buoyancy, friction) and simple machines (levers,
pulleys) within one coherent environment and interaction model. Our sandbox
targets this intersection. We use pseudo-haptic gain and visual-lag modulation as
a unifying mechanism that makes an unrenderable physical quantity perceptible
across a family of controlled, single-variable scenarios, and we extend the same
mechanism from holding a heavier object to feeling a loaded lever push back.
Crucially, the cue tracks \emph{mass} for held objects and \emph{load} for the
lever, the quantities that actually determine weight and effort, so that density
is taught where it belongs, through the volume difference between equal-mass
blocks and through float-or-sink behavior, rather than being conflated with
heaviness. Because our contribution is an interaction technique for making a
physical quantity felt, we evaluate it at two levels: a system level, checking
that the cues behave monotonically with the underlying physics and that each
scenario matches its closed-form prediction; and a human level, a within-subjects
study testing whether the weight cue changes how accurately learners predict
outcomes that depend on mass, density, and load.
 in main.tex in place of the current
%% \section{Related Work}.
%% ============================================================================

\section{Related Work}
Our sandbox sits at the intersection of four lines of research: immersive
virtual reality (VR) for science learning, the difficulty of rendering weight and
force on commodity hardware, the pseudo-haptic techniques that stand in for that
missing sensation, and the interaction-design foundations that make direct
manipulation legible. We review each in turn and close by stating what
distinguishes our approach.

\subsection{Immersive VR for Science Learning}
The pedagogical case for immersive VR rests less on visual fidelity than on the
sense of being present. Slater and S{\'a}nchez-Vives frame this through the place
illusion (the feeling of being in the virtual space) and the plausibility
illusion (the feeling that events are really happening), which together lead
users to respond to virtual events as if they were real~\cite{slater2016enhancing}.
For a physics sandbox this matters directly. If a learner treats a virtual block
and lever as real objects, then manipulating them can carry the intuitive weight
that a 2D simulation cannot.

Empirical work in physics education adds a consistent caveat: immersion helps,
but only when the learning experience is designed for it. A recent meta-analysis
by Wu et al.\ across many head-mounted display studies finds that immersive VR
yields a positive but moderate overall effect on learning performance, and that
this effect is strongly moderated by factors such as subject domain, learner
control, and how the experience is designed~\cite{wu2020effectiveness}.
Tsivitanidou, Georgiou, and Ioannou report on an inquiry-based immersive VR
physics activity from two angles: the learning experience they observe conceptual
gains alongside an interaction with learners' attitudinal
profiles~\cite{tsivitanidou2021inquiry}, and the design of that same experience,
which they document as a worked case of learning-experience
design~\cite{georgiou2021led}. We read this pair as evidence that immersive VR
physics learning is feasible and can produce conceptual gains; it establishes
that a VR physics activity can work, though it does not itself isolate felt weight
or the effect of a haptic cue. One common way to scaffold such experiences is an
embedded pedagogical agent. Makransky et al.\ studied one such agent and found a
gender-matching effect: a virtual instructor aided learning most when its
apparent gender matched the learner, a specific interaction rather than a general
endorsement of pedagogical agents~\cite{makransky2019gender}. Taken together,
these studies establish that immersive physics learning is promising but
design-sensitive. The value comes from what the learner is asked to do and how
the task is scaffolded, not from the headset itself. We take a deliberately
different scaffolding stance. Rather than an in-scene agent, we scaffold through
scene and task design, using tightly scoped scenarios that isolate a single
variable so that the cause of each outcome is unambiguous.

\subsection{The Missing Channel: Rendering Weight and Force}
The central obstacle for a physics-of-materials experience is that the very
quantity the lessons depend on, namely weight and the force needed to move a
loaded mechanism, is the one commodity that VR cannot deliver. A handheld
controller weighs the same whatever it represents. Wang et al.\ survey the state
of haptic display for VR and make clear that genuine, general-purpose force
feedback remains far from consumer reach, constrained by actuation, workspace,
and grounding requirements that headset-and-controller setups do not
meet~\cite{wang2019haptic}. Their survey frames the design problem for us
precisely. If we want a denser object to feel heavier on a Quest 2, we cannot
wait for the hardware to render the force, and must instead substitute perception
for physics.

\subsection{Pseudo-Haptics: Substituting Perception for Physics}
Pseudo-haptics addresses this gap by exploiting the dominance of vision over
proprioception. Manipulating the visual consequences of a motion, such as the
control-to-display gain or a small lag between hand and object, induces a felt
sense of weight, stiffness, or resistance without any force actuator. The
foundational demonstration comes from L{\'e}cuyer et al., who showed that
altering the visual response of an isometric input device makes users perceive
force and stiffness that no actuator produces, establishing that a
control-to-display manipulation alone can stand in for force
feedback~\cite{lecuyer2000pseudo}. Ujitoko and Ban survey the design space and
applications of pseudo-haptics, cataloguing how such visual manipulations
reliably produce haptic sensations and how they have been
applied~\cite{ujitoko2021pseudo}. Their taxonomy is the methodological backbone
for how we map mass and mechanical load onto gain and visual lag.

Two results ground our specific cue, the felt \emph{weight} of a held object.
Dominjon et al.\ manipulated the control/display ratio of a manipulated object in
a virtual environment and showed that this ratio directly changes its perceived
mass, so that the same physical motion can be made to feel heavier or lighter by
scaling its visual displacement~\cite{dominjon2005influence}. Rietzler et al.\
carried the same idea to consumer VR: by introducing perceivable offsets between
the tracked hand and the rendered object, they elicited a reliable sensation of
weight on a commodity headset with no force
hardware~\cite{rietzler2018breaking}. Together these establish that a gain or
tracking-offset manipulation can render felt weight on exactly the class of
device we target, and they fix the design variable for us: the cue must track an
object's \emph{mass}, the quantity that determines weight, rather than its
density. These techniques are not limited to lifting a single object. The same
principle of modulating the visual outcome of an effort extends naturally to the
resistance of operating a machine, such as a lever that grows harder to move as
its load increases.

\subsection{Foundations of 3D Interaction}
Whatever the illusion, the learner must first be able to reach out, grab, and
place objects fluidly, or the physics never gets probed. We ground our
interaction design in Bowman et al.'s canonical treatment of 3D user interfaces,
which formalizes selection, manipulation, and travel as the core tasks of spatial
interaction and catalogues the techniques for each~\cite{bowman2004}. This
foundation informs the grab-and-place model at the heart of the sandbox and keeps
our contribution, the pseudo-haptic property mapping, layered on top of
well-understood manipulation techniques rather than reinventing them.

\subsection{What Is Different About Our Approach}
Two gaps emerge from this literature. First, immersive physics studies
consistently deliver visual and interactive experiences but rarely give learners
a genuine felt contrast in weight between objects. A student may see that two
blocks differ, but does not feel one resist more than the other. Second,
pseudo-haptic studies convincingly render the felt weight of an isolated
object~\cite{dominjon2005influence, rietzler2018breaking}, but seldom span both
material properties (density, buoyancy, friction) and simple machines (levers,
pulleys) within one coherent environment and interaction model. Our sandbox
targets this intersection. We use pseudo-haptic gain and visual-lag modulation as
a unifying mechanism that makes an unrenderable physical quantity perceptible
across a family of controlled, single-variable scenarios, and we extend the same
mechanism from holding a heavier object to feeling a loaded lever push back.
Crucially, the cue tracks \emph{mass} for held objects and \emph{load} for the
lever, the quantities that actually determine weight and effort, so that density
is taught where it belongs, through the volume difference between equal-mass
blocks and through float-or-sink behavior, rather than being conflated with
heaviness. Because our contribution is an interaction technique for making a
physical quantity felt, we evaluate it at two levels: a system level, checking
that the cues behave monotonically with the underlying physics and that each
scenario matches its closed-form prediction; and a human level, a within-subjects
study testing whether the weight cue changes how accurately learners predict
outcomes that depend on mass, density, and load.
